\documentclass[11pt]{article}
\usepackage{amsmath,amssymb,amsthm,mathtools}
\usepackage[margin=1in]{geometry}
\usepackage{hyperref}

\newtheorem{theorem}{Theorem}
\newtheorem{lemma}[theorem]{Lemma}
\newtheorem{proposition}[theorem]{Proposition}
\newtheorem{definition}[theorem]{Definition}
\newtheorem{remark}[theorem]{Remark}
\theoremstyle{definition}
\newtheorem{verified}[theorem]{Computational Result}

\newcommand{\Z}{\mathbb{Z}}
\newcommand{\Q}{\mathbb{Q}}
\newcommand{\K}{K}
\newcommand{\OK}{\mathcal{O}_K}
\newcommand{\zetaK}{\zeta_{K}}
\newcommand{\Icos}{\mathcal{I}}
\renewcommand{\Theta}{\vartheta}

\title{Geometric Realizations of $L$-Functions on the 600-Cell}
\author{(VFD programme --- working draft)}
\date{2026}

\begin{document}
\maketitle

\begin{abstract}
We record explicit correspondences between sub-objects of the $600$-cell / icosian
ring and classical $L$-functions, together with their computational verification.
Three correspondences are exhibited: (i) the icosian theta series over
$K=\Q(\sqrt5)$ realizes $\zetaK(s)\,\zetaK(s-1)$; (ii) the rational $24$-cell
sub-polytope realizes $\zeta(s)\,\zeta(s-1)$, whose zeros on the critical line are
exactly the nontrivial zeros of the Riemann zeta function; and (iii) the
Galois ``twin'' (the $\sqrt5$-shell of $96$ vertices) carries the quadratic
Dirichlet $L$-function $L(s,\chi_5)$, in accordance with the classical
factorization $\zetaK=\zeta\cdot L(\chi_5)$. We state the scope precisely: these are
\emph{realizations and verifications of standard $L$-functions} attached to known
lattices and number fields. \textbf{No claim is made about the Riemann Hypothesis,
and no claim is made of a bridge to any physical model.}
\end{abstract}

\section{Introduction and scope}
The $600$-cell is the regular $4$-polytope whose $120$ vertices are the unit
icosians, i.e.\ the binary icosahedral group $2I\subset\mathbb{H}$. Its vertex set
carries rich arithmetic: the icosian ring is a maximal order in a quaternion algebra
over $K=\Q(\sqrt5)$, and the associated lattice is a copy of $E_8$
\cite{ConwaySloane}. This note isolates three \emph{exact} arithmetic objects
attached to $600$-cell sub-structures and verifies, numerically, which $L$-functions
they realize.

\paragraph{What is claimed.} Only the correspondences in
\S\ref{sec:icosian}--\S\ref{sec:twin} and their numerical verification
(\S\ref{sec:verify}). All number-theoretic inputs (Dedekind factorization,
functional equations) are standard and cited.

\paragraph{What is not claimed.} Nothing about the location of the zeros of
$\zeta$ (the Riemann Hypothesis); no positivity, operator, or Frobenius statement;
no connection to any cosmological or physical model. These boundaries are
deliberate and load-bearing.

\section{Preliminaries}
Let $K=\Q(\sqrt5)$, $\varphi=\tfrac{1+\sqrt5}{2}$, $\OK=\Z[\varphi]$. The Dedekind
zeta function $\zetaK(s)$ factors classically as
\begin{equation}\label{eq:dedekind}
  \zetaK(s)=\zeta(s)\,L(s,\chi_5),
\end{equation}
where $\chi_5$ is the (even, real) quadratic Dirichlet character modulo $5$
\cite{Neukirch}. Both factors satisfy functional equations; $\zeta$ has a simple
pole at $s=1$, $L(\cdot,\chi_5)$ is entire.

The icosian ring $\Icos$ is the $\Z[\varphi]$-order generated by the $120$ unit
icosians (the $600$-cell vertices). Its $24$-element rational sub-structure is the
ring of Hurwitz units (the $24$-cell / binary tetrahedral group $2T$), defined over
$K=\Q$. The remaining $96$ vertices form the $\varphi$-shell; the coordinate Galois
involution $\sigma:\sqrt5\mapsto-\sqrt5$ fixes the $24$-cell pointwise and permutes
the $96$ $\varphi$-vertices to the odd-permutation twin.

For a (positive-definite) lattice $\Lambda$ we write $\Theta_\Lambda(\tau)=\sum_{x
\in\Lambda}q^{\langle x,x\rangle/2}$, $q=e^{2\pi i\tau}$, for its classical theta
series, and $L(\Theta_\Lambda,s)$ for the standard $L$-function of the associated
modular form. In \S\ref{sec:icosian}, where the relevant form lives over
$K=\Q(\sqrt5)$, $\Theta$ denotes instead the corresponding \emph{Hilbert} theta series
(over the two infinite places of $K$); we make this explicit there.

\section{The icosian theta series and \texorpdfstring{$\zetaK(s)\zetaK(s-1)$}{zetaK(s)zetaK(s-1)}}\label{sec:icosian}
The clean identity below rests on one arithmetic input---that the icosian order has a
single class in its genus. We isolate this input as a lemma and pin it to the
literature, since it is the load-bearing hypothesis of the section.

\begin{lemma}[class number one]\label{lem:h1}
Let $B$ be the definite quaternion algebra over $K=\Q(\sqrt5)$ ramified exactly at the
two infinite places and unramified at all finite places, and let $\Icos\subset B$ be
the icosian maximal order (the $\Z[\varphi]$-order on the $120$ unit icosians). Then
$\Icos$ has \emph{class number one}: a single left-ideal class.
\end{lemma}
\begin{proof}[Reference]
The pair $(B,\Icos)$ is the maximal-order case with invariants
$(n,d_F,D,N)=(2,5,1,1)$ in the Kirschmer--Voight classification of definite quaternion
orders of class number one over totally real fields
\cite[Table 8.2]{KirschmerVoight} (see also Vign\'eras \cite{Vigneras}); it is the
classical class-number-one example over $\Q(\sqrt5)$. This is the left-ideal class
number of the order, the quantity entering the Siegel--Weil average below. (The
underlying $\Z$-lattice of $\Icos$ is $E_8$; we note this only as context---$E_8$'s
uniqueness as an even unimodular $\Z$-lattice is a statement about $\Z$-lattice genera,
not about quaternion ideal classes, and is \emph{not} used here.)
\end{proof}

\begin{proposition}\label{res:icosian}
Let $\Theta_{\Icos}$ be the theta series of the icosian order over $K=\Q(\sqrt5)$.
Then
\[
   L\bigl(\Theta_{\Icos},s\bigr)=\zetaK(s)\,\zetaK(s-1).
\]
\end{proposition}
\begin{proof}[Proof sketch]
We take $\Theta_{\Icos}$ to be the Hilbert theta series over $K$ of the quaternionic
norm form on $\Icos$ (weight two, the level determined by the discriminant
$(D,N)=(1,1)$), normalized so that its first nonzero Fourier coefficient is $1$; and
$L(\Theta_{\Icos},s)$ its standard (Hecke) $L$-function. By the Siegel--Weil formula
\cite{Siegel,Weil1965}, the genus-average of such a theta series is the corresponding
Eisenstein series. By Lemma~\ref{lem:h1} the order has a single ideal class, so the
average reduces to the single term $\Theta_{\Icos}$; hence $\Theta_{\Icos}$ \emph{is}
that Eisenstein series, with no cuspidal contribution. The standard $L$-function of the
relevant weight-two Hilbert Eisenstein series over $K$ is the rank-two product
$\zetaK(s)\zetaK(s-1)$ (constant-term / Hecke-eigenvalue computation; see
\cite{Eichler,HMF}). We present this as a \emph{sketch}: the precise Eisenstein
normalization and local factors are standard but not reproduced here, and the
identity's force in this note rests jointly on the sketch and the numerical
corroboration below.
\end{proof}
\begin{remark}
Class number one is what makes the identity clean: were the order to have $h>1$ ideal
classes, $\Theta_{\Icos}$ would decompose as the Eisenstein series \emph{plus} a cusp
form, and $L(\Theta_{\Icos},s)$ would be a sum carrying additional cuspidal
$L$-contributions rather than the bare product $\zetaK(s)\zetaK(s-1)$.
Lemma~\ref{lem:h1} excludes this.
\end{remark}
\begin{remark}[numerical corroboration; provenance]
The identity was corroborated numerically in the companion icosian/Brandt computation
of this programme (the \texttt{icosian-triad} work), where the representation numbers
$r(\pi)$ were obtained by exact short-vector enumeration over $\Z[\varphi]$ and matched,
out of sample, against the Dirichlet coefficients of $\zetaK(s)\zetaK(s-1)$. We flag
provenance honestly: the verification-engine scripts shipped \emph{with this bundle}
(\S\ref{sec:verify}, \texttt{vfd\_math\_engine}) illustrate the gate by declaring the
predicted zero lists directly; they do \emph{not} re-derive $r(\pi)$ from the lattice.
The geometry$\to$arithmetic enumeration is the separate computation just cited.
\end{remark}

\section{The rational \texorpdfstring{$24$}{24}-cell and \texorpdfstring{$\zeta(s)\zeta(s-1)$}{zeta(s)zeta(s-1)}}\label{sec:24cell}
Applying the same construction to the rational sub-order (the $24$-cell, $K=\Q$)
yields the $\Q$-analogue of Result~\ref{res:icosian}.
\begin{verified}\label{res:24cell}
The $24$-cell (Hurwitz) theta series realizes
\[
   L(\Theta_{24},s)=\zeta(s)\,\zeta(s-1).
\]
\end{verified}
\begin{proposition}\label{prop:online}
On the critical line $\mathrm{Re}(s)=\tfrac12$, the zeros of $\zeta(s)\zeta(s-1)$ are
exactly those nontrivial zeros of $\zeta$ that lie on that line. (Equivalently: the
factor $\zeta(s-1)$ contributes none of them.)
\end{proposition}
\begin{proof}
A nontrivial zero of $\zeta(s-1)$ occurs at $s-1=\rho$, i.e.\ $s=1+\rho$; since every
nontrivial zero $\rho$ has $\mathrm{Re}(\rho)\in(0,1)$, such $s$ has
$\mathrm{Re}(s)\in(1,2)$, never $\tfrac12$. The trivial zeros of $\zeta(s-1)$ lie at
$s=1-2n$ ($n\ge1$), also off the line. Hence on $\mathrm{Re}(s)=\tfrac12$ the factor
$\zeta(s-1)$ contributes no zeros, and the on-line zeros of the product are precisely
the on-line zeros of $\zeta(s)$. This is a statement about \emph{which} zeros lie on the
line, and asserts nothing about whether all nontrivial zeros of $\zeta$ do
(see \S\ref{sec:scope}).
\end{proof}
Thus the $24$-cell --- the $\sqrt5$-free, Galois-fixed sub-polytope --- carries the
Riemann zeros on its critical line, with the second factor displaced to
$\mathrm{Re}(s)=\tfrac32$. (This is a statement about \emph{which} zeros lie on the
line, not about whether all zeros of $\zeta$ do; see \S\ref{sec:scope}.)

\section{The Galois twin and \texorpdfstring{$L(s,\chi_5)$}{L(s,chi5)}}\label{sec:twin}
The factorization \eqref{eq:dedekind} predicts that the ``$\sqrt5$ part'' of the
icosian structure carries $L(\chi_5)$. Geometrically, this is the $\varphi$-shell.
\begin{verified}\label{res:twin}
The $96$-vertex $\varphi$-shell (the Galois-sign component) realizes the Dirichlet
$L$-function $L(s,\chi_5)$: its on-line zeros coincide (to tested precision) with the
zeros of $L(s,\chi_5)$; no coincidence with a Riemann zero was observed in the tested
range.
\end{verified}
Together, Results~\ref{res:icosian}--\ref{res:twin} are mutually consistent via
\eqref{eq:dedekind}: the full icosian object carries $\zetaK(s)\,\zetaK(s-1)$
$=\zeta(s)\,L(s,\chi_5)\,\zeta(s-1)\,L(s-1,\chi_5)$; its Galois-fixed part carries the
$\zeta$-factor (the $24$-cell), and its Galois-sign part carries the $L(s,\chi_5)$
factor (the $\varphi$-shell).

\section{Verification}\label{sec:verify}
The correspondences were checked computationally. For each claimed $L$-function we
compared its on-line zeros against the target's zeros on four invariants: (a) first-$k$
zero positions; (b) zero density (the Weyl/Riemann--von~Mangoldt rate); (c) the
nearest-neighbour spacing distribution (``\,fp'' below is the fraction of normalized
gaps below $0.3$); (d) the number variance $\Sigma^2(\ell)$. The geometry$\to$arithmetic
input---the icosian representation numbers $r(\pi)$ by exact short-vector enumeration
over $\Z[\varphi]$, matched out of sample against $\zetaK(s)\zetaK(s-1)$---was performed
in the companion icosian/Brandt computation (cited in \S\ref{sec:icosian}). The engine
scripts in \emph{this} bundle exercise the verification gate on declared zero lists and
are not the enumeration; the figures below are reproduced from the cited computation.

\begin{center}\small
\begin{tabular}{lll}
\hline
correspondence (on-line content) & first on-line zeros & spacing fp ($<0.3$)\\
\hline
$24$-cell $\to \zeta(s)\zeta(s-1)$ & $14.13,\,21.02,\,25.01,\,30.42$ & $0.000$\\
$\varphi$-shell $\to L(s,\chi_5)$ & $6.65,\,9.83,\,11.96,\,16.03$ & $0.000$\\
icosian $\to \zetaK(s)\zetaK(s-1)$ & $6.65,\,9.83,\,11.96,\,14.13$ & $0.077$\\
\hline
\end{tabular}
\end{center}
\noindent The single-$L$-function rows are consistent with GUE nearest-neighbour
repulsion (fp $\approx 0$); the rank-four icosian row shows the larger value expected
for a \emph{superposition} of independent spectra
($\zetaK(s)\zetaK(s-1)=\zeta(s)\,L(s,\chi_5)\,\zeta(s-1)\,L(s-1,\chi_5)$), consistent
with \eqref{eq:dedekind}. These are diagnostic comparisons against the named target
$L$-functions, not independent evidence for the geometric attachment.

\section{Scope and what is not claimed}\label{sec:scope}
We restate the boundary. Propositions~\ref{prop:online} and
Results~\ref{res:icosian}--\ref{res:twin} describe \emph{which} $L$-functions are
realized by which $600$-cell sub-objects, and which of their zeros lie on the
critical line. They are \textbf{silent on the Riemann Hypothesis}: the claim
``the $24$-cell carries the Riemann zeros'' refers to the zeros that \emph{do} lie
on $\mathrm{Re}(s)=\tfrac12$, and asserts nothing about whether \emph{all}
nontrivial zeros of $\zeta$ lie there. No positivity, self-adjointness, or
arithmetic-surface statement is made or implied. A companion note records, with
explicit certificates, that the natural ``substrate $\leftrightarrow\zeta$''
correspondence is \emph{negative} --- the icosian object realizes
$\zetaK(s)\zetaK(s-1)=\zeta(s)L(s,\chi_5)\zeta(s-1)L(s-1,\chi_5)$, which contains
$\zeta$ only as a factor and does not isolate it.

\begin{thebibliography}{9}
\bibitem{ConwaySloane} J.~H.~Conway, N.~J.~A.~Sloane,
  \emph{Sphere Packings, Lattices and Groups}, Springer.
\bibitem{Neukirch} J.~Neukirch, \emph{Algebraic Number Theory}, Springer.
\bibitem{Eichler} M.~Eichler, \emph{The basis problem for modular forms and the
  traces of the Hecke operators}, Lecture Notes in Math.
\bibitem{HMF} E.~Freitag, \emph{Hilbert Modular Forms}, Springer.
\bibitem{Siegel} C.~L.~Siegel, \emph{\"Uber die analytische Theorie der
  quadratischen Formen}, Ann.\ of Math.\ (1935).
\bibitem{Weil1965} A.~Weil, \emph{Sur la formule de Siegel dans la th\'eorie des
  groupes classiques}, Acta Math.\ (1965).
\bibitem{KirschmerVoight} M.~Kirschmer, J.~Voight, \emph{Algorithmic enumeration of
  ideal classes for quaternion orders}, SIAM J.\ Comput.\ \textbf{39} (2010).
\bibitem{Vigneras} M.-F.~Vign\'eras, \emph{Arithm\'etique des alg\`ebres de
  quaternions}, Lecture Notes in Math.\ \textbf{800}, Springer (1980).
\end{thebibliography}

\end{document}
